\section{Algebraic Spine and Boundary Theorem}
\label{sec:theorem}

The theorem below states the finite boundary result in public notation.
Each finite assertion is tied to a source artifact in
\(\texttt{paper/SOURCE\_MAP.md}\); the proof uses those artifacts as certified
finite inputs rather than as hidden mathematical assumptions.

\begin{certprop}[Line-xor and degree-bound certificate]
\label{cert:line-xor-degree}
In value-minus-one labels, every row and every column of each certified normal
order-four magic square has xor zero.  This is the local Nimm0/Lemma B input.
The manuscript treats the Step4 finite derivation as the proof source; the
classical order-four literature supplies context rather than an unverified
quoted theorem~\cite{ollerenshawBondi,wardMagicVectorSpaces}.  With the
row/column line-xor input, the degree-bound argument gives
\[
  \deg L\le 2.
\]
\end{certprop}

\begin{proof}
Write each output bit of \(L(r_1,r_0,c_1,c_0)\) in algebraic normal form over
\(\F_2\).  For a monomial \(r^{S_r}c^{S_c}\), xoring over a fixed row leaves it
alive exactly when \(S_c=\{c_1,c_0\}\).  Hence the condition that every row xors
to zero kills precisely the monomials divisible by \(c_1c_0\).  Similarly,
column-xor zero kills precisely the monomials divisible by \(r_1r_0\).  Every
degree-three or degree-four monomial in the four variables is divisible by one
of these two pairs, so no monomial of degree at least three remains.  The
row/column line-xor premise itself is the Step4 finite row-pair derivation from
the full row, column, and diagonal magic constraints.
\end{proof}

\begin{lemma}[Quadratic error form]
\label{lem:quadratic-error-form}
Let \(A\) be the affine interpolation of \(L\) on
\(\{0,e_1,e_2,e_3,e_4\}\), and set \(E=L+A\).  For the certified normal
magic-square records, the pure row-pair and pure column-pair quadratic
coefficients vanish.  Hence the error has bilinear form
\[
  E(r,c)=B_{11}r_1c_1+B_{10}r_1c_0+B_{01}r_0c_1+B_{00}r_0c_0.
\]
\end{lemma}

\begin{proof}
The affine interpolation \(A\) absorbs the constant and linear parts of \(L\).
By Proposition~\ref{cert:line-xor-degree}, no terms of degree at least three
occur.  The row-xor condition kills the pure column-pair term \(c_1c_0\), while
the column-xor condition kills the pure row-pair term \(r_1r_0\).  The only
quadratic terms left are the four mixed row/column monomials displayed above.
\end{proof}

\begin{theorem}[Defect dichotomy / Lemma C]
\label{thm:defect-dichotomy}
For the degree-at-most-two bijections of \(\F_2^4\) covered by the certificate,
the bilinear coefficient matrix has the certified one-dependent-side property.
Consequently the error-direction span \(\Derr\) has dimension at most two and
\[
  (\delta,\dim\Derr)\in\{(0,0),(4,1),(6,2)\}.
\]
This is the first-principles structural part of the affine-defect spine.
\end{theorem}

\begin{proof}
Let \(B=(B_{ij})\) be the \(2\times2\) matrix of mixed bilinear coefficients.
The property that a bilinear matrix is completable to a degree-at-most-two
bijection \(x\mapsto A_0x+E(x)\), and the property that one side of \(B\) is
dependent, are invariant under the codomain action of \(GL(4,2)\).  The
Step5+ check therefore reduces the obstruction to a symmetry-complete finite
set: all \(4096\) matrices with entries in a fixed three-dimensional subspace,
plus one representative for the four-dimensional span case.  In that reduced
universe every completable matrix has one dependent side, and the
four-dimensional representative is not completable.  Thus a quadratic
bijection has \(\dim\Derr\le2\) and a dependent side.

With one dependent side, \(E\) factors as a nonzero scalar linear form on one
\(\F_2^2\) coordinate side times a vector-valued linear map on the other side.
If \(d=\dim\Derr\), the scalar form is nonzero on two points and the vector map
is nonzero on \(4-2^{2-d}\) points.  Hence
\[
  \delta=2(4-2^{2-d}),
\]
which gives exactly \((0,0),(4,1),(6,2)\).
\end{proof}

\begin{lemma}[Error-fiber parity bridge]
\label{lem:error-fiber-parity}
Let \(M\) be the selected mask.  If \(\delta=4\), then
\(\Derr=\{0,u\}\) and \(E=u\) on a four-cell support \(\Sigma\), so
\[
  \bigoplus_{c\in M} L(c)=u\cdot(|M\cap\Sigma|\bmod 2).
\]
If \(\delta=6\), the three nonzero fibers
\(\Sigma_u,\Sigma_v,\Sigma_{u+v}\) contribute
\[
  \bigoplus_{c\in M}L(c)=
  (|M\cap\Sigma_u|\bmod2)u+
  (|M\cap\Sigma_v|\bmod2)v+
  (|M\cap\Sigma_{u+v}|\bmod2)(u+v).
\]
Thus selected-label affineness is controlled by error-fiber parity.
\end{lemma}

\begin{proof}
The selected masks in this endpoint atlas are affine planes in the cell space,
so the affine part cancels:
\[
  \bigoplus_{c\in M}A(c)=0.
\]
Thus \(\bigoplus_{c\in M}L(c)=\bigoplus_{c\in M}E(c)\).  The factorization in
Theorem~\ref{thm:defect-dichotomy} gives the fiber shapes: for \(\delta=4\) a
single nonzero error value \(u\) occurs on a four-cell support, and for
\(\delta=6\) the three nonzero values \(u,v,u+v\) occur on three two-cell
fibers.  Grouping the xor of \(E\) over \(M\) by these fibers gives exactly the
two displayed parity formulas.
\end{proof}

\begin{certprop}[Free \(\TK\)-transport]
\label{cert:tk-transport}
The finite transport certificate proves a free global \(\TK\cong V_4\) action on
the certified atlas.  It preserves endpoint, class, defect, and selected
signature.  Therefore endpoint \(24\) has \(236=59\cdot4\), with quotient split
\[
  59=36+8+15.
\]
\end{certprop}

\begin{certprop}[B2 selected-affine boundary]
\label{cert:b2-boundary}
Inside endpoint \(24\), B2 is the selected-affine, cell-value-non-affine
boundary.  It has \(32\) records and \(8\) free \(\TK\)-orbits.  In the
\(\delta=4\) fiber picture it is the even side of the constant-error support
parity:
\[
  |M\cap\Sigma|\equiv0\pmod2.
\]
The endpoint package verifies this selected-affine, cell-value-non-affine
finite selector and closes the \(32/8\) count inside the saved endpoint atlas.
This is a structural parity selector plus a finite endpoint package, not a
complete conceptual B2 classification theorem.
\end{certprop}

\begin{certprop}[B3 support-action matrix]
\label{cert:b3-support-action}
Inside endpoint \(24\), B3 is the selected-non-affine boundary.  It has
\(60\) records and \(15\) free \(\TK\)-orbits.  The support-action matrix has
eight non-affine selected signatures against eight support-action profiles,
hence \(64\) candidate cells.  The finite bilinear transport package verifies
exactly \(15\) realized cells and \(49\) absent cells, with no false positives
or false negatives.  The intrinsic column-classification certificate refines
the column side: all eight profile columns are intrinsic to the selected
signature and the bilinear error matrix \(B\), through \(\delta\), the zero
pattern of \(B\), and \(\Derr=\operatorname{span}(B)\).  The defect-6
realizability-obstruction certificate identifies the remaining realization side
as an integer-magic residue: at defect \(6\), the order-four magic-square
constraints realize \(13\) of the \(35\) two-dimensional error spaces, although
all \(35\) are \(\F_2\)-completable and the realized \(13\)-set has trivial
\(GL(4,2)\) stabilizer.  Thus the column classification is intrinsic, while
the \(15/49\) realization cut remains a finite arithmetic certificate.  The
NF1/NF2/NF3 templates are structural clause shapes; the thirteen profile-level
instances are finite-certified.
\end{certprop}

\begin{certprop}[B3 count collapse]
\label{cert:b3-count-collapse}
The quotient count \(15\) is not an independent count once freeness is known:
\[
  15=60/4.
\]
The finite residual is the row-degree pattern
\[
  1,1,1,2,2,1,4,3
\]
and the endpoint quotient kernel
\[
  \delta=0,4,6\mapsto36,16,7.
\]
This is a finite debt reduction, not a symbolic proof of the count \(60\).  A
non-enumerative derivation of that count remains outside the claim.
\end{certprop}

\begin{certprop}[B1 scoped affine-plane absence]
\label{cert:b1-absence}
B1 is the scoped endpoint \(23/27\) absence statement for the affine
selected-plane layer.  In that layer every selected four-set has selected-label
xor zero.  The endpoint \(23/27\) signatures under discussion have nonzero
selected-label xor, and therefore cannot occur in the affine selected-plane
layer.  This is not a full endpoint \(23/27\) classification.
\end{certprop}

\begin{theorem}[Affine-defect boundary atlas, certified endpoint-24 form]
\label{thm:b123-affine-defect-boundary}
Inside the certified endpoint-24 order-four magic-square atlas, consider square-mask pairs in
the bounded one-incidence model.  Quotient statements are taken modulo the
verified \(D_4\) essential representatives and the verified free \(\TK\)
action.  Then:
\begin{enumerate}[leftmargin=*]
\item endpoint \(24\) has \(236\) records and \(59\) free \(\TK\)-orbits;
\item the affine defect satisfies \(\delta\in\{0,4,6\}\), with record split
  \(144/64/28\);
\item the exact affine core has \(144\) records and \(36\) \(\TK\)-orbits;
\item B2 has \(32\) records and \(8\) \(\TK\)-orbits;
\item B3 has \(60\) records and \(15\) \(\TK\)-orbits.  Its support-action
  matrix has \(64\) candidate cells, exactly \(15\) realized cells, and
  exactly \(49\) absent cells;
\item B1 is the scoped endpoint \(23/27\) absence statement from the affine
  selected-plane layer, not a full endpoint \(23/27\) classification.
\end{enumerate}
\end{theorem}

\begin{proof}
Proposition~\ref{cert:line-xor-degree} gives \(\deg L\le2\) from row/column
line xor zero, and Lemma~\ref{lem:quadratic-error-form} writes \(E=L+A\) as a
bilinear error on row/column coordinates.  Theorem~\ref{thm:defect-dichotomy}
gives \(\delta\in\{0,4,6\}\) with
\((\delta,\dim\Derr)=(0,0),(4,1),(6,2)\), and the error-fiber parity bridge
identifies the selected-label xor from the selected mask's intersections with
the error fibers.

The free \(\TK\)-transport proposition gives \(236=59\cdot4\) on endpoint
\(24\).  The exact affine core is the selected-affine, \(\delta=0\) cell and
has \(144\) records, hence \(36\) free \(\TK\)-orbits.  The B2 proposition gives
the selected-affine boundary as \(32\) records and \(8\) orbits.  The B3
support-action proposition gives the selected-non-affine boundary as \(60\)
records and \(15\) orbits, with the \(15/49\) support-action matrix.  The
count-collapse proposition records the sharper finite residual without turning
it into a symbolic count proof.  Finally, the B1 proposition gives the scoped
endpoint \(23/27\) absence from the affine selected-plane layer.  These blocks
are exactly the assertions of the theorem.
\end{proof}

\begin{certprop}[Full endpoint aggregate controls]
\label{cert:full-endpoint-control}
The public aggregate control artifact verifies the full
\(880\cdot 8=7040\) square--mask universe used as context for this paper.  For
every square--mask pair with label map \(L\) and selected mask \(M\),
\[
  \operatorname{end}(L,M)=34-\min_{c\in M}L(c)
  =35-\min\{\text{selected values}\}.
\]
The endpoint distribution is
\[
\begin{array}{c|rrrrrrrrrrrrr}
 e&22&23&24&25&26&27&28&29&30&31&32&33&34\\\hline
 \#&176&44&236&268&676&108&328&376&772&420&884&992&1760.
\end{array}
\]
The same artifact gives the aggregate \(T/K\) endpoint-orbit counts
\[
44,11,59,67,169,27,82,94,193,105,221,248,440
\]
for endpoints \(22,23,\ldots,34\).  Endpoints \(23\) and \(27\) have no
possible selected-affine signature, while endpoint \(22\) is a positive
selected-affine control: all its \(176\) aggregate records are selected-affine.
Endpoint \(24\) is therefore the first anchor treated here where the
selected-affine and selected-non-affine layers both occur, with the certified
split \(176+60=236\) and affine-defect distribution \(144,64,28\).

This is a finite aggregate control statement.  It is not a classification of
every endpoint and it does not widen
Theorem~\ref{thm:b123-affine-defect-boundary} beyond endpoint \(24\).
\end{certprop}
